% Partie : set-theory
% Chapitre : ordinals
% Section : iso

\documentclass[../../../include/open-logic-section]{subfiles}

\begin{document}

\olfileid[fr]{sth}{ordinals}{iso}
\olsection{Isomorphismes d'ordre}

Pour expliquer \emph{à quel point} la notion de bon ordre est robuste,
commençons par introduire une méthode de comparaison des bons ordres.

\begin{defn}
Un \emph{bon ordre} est un couple $\tuple{A, <}$ tel que $<$ définit
un bon ordre sur $A$. Les bons ordres $\tuple{A, <}$ et $\tuple{B,
\lessdot}$ sont \emph{isomorphes en tant qu'ordres} si et seulement
s'il existe !!a{bijection} $f \colon A \to B$ telle que : $x < y$ si
et seulement si $f(x) \lessdot f(y)$. Dans ce cas, nous écrivons
$\ordeq{\tuple{A, <}}{\tuple{B, \lessdot}}$ et disons que $f$ est un
\emph{isomorphisme d'ordre}.
\end{defn}
\noindent
Dans la suite, par souci de brièveté, nous parlerons simplement
d'«~isomorphismes~» plutôt que d'«~isomorphismes d'ordre~».
Intuitivement, les isomorphismes sont des !!{bijection}s qui préservent
la structure. Voici quelques faits simples à leur sujet.

\begin{lem}\ollabel{isoscompose}
La composée de deux isomorphismes est un isomorphisme : si $f \colon A
\to B$ et $g \colon B \to C$ sont des isomorphismes, alors
$(g \circ f) \colon A \to C$ est un isomorphisme.
\end{lem}
\begin{prob}
	Démontrez le \olref[sth][ordinals][iso]{isoscompose}.
\end{prob}
\begin{proof}
La démonstration est laissée en exercice.
\end{proof}
\begin{cor}\ollabel{ordisoisequiv}
	$\ordeq{X}{Y}$ est une relation d'équivalence.
\end{cor}
\begin{prop}\ollabel{ordisounique}
Si $\tuple{A, <}$ et $\tuple{B, \lessdot}$ sont des bons ordres
isomorphes, l'isomorphisme entre eux est unique.
\end{prop}

\begin{proof}
Soient $f$ et $g$ deux isomorphismes de $A$ sur $B$. Nous allons
démontrer le résultat par récurrence, en utilisant
la \olref[wo]{propwoinduction}.
Fixons $a\in A$ et supposons, comme hypothèse de récurrence, que
$(\forall b < a)f(b) = g(b)$. Fixons aussi $x \in B$.

Si $x \lessdot f(a)$, alors $f^{-1}(x) < a$ et donc
$g(f^{-1}(x)) \lessdot g(a)$, puisque $f$ et $g$ sont des
isomorphismes. Mais, comme $f^{-1}(x) < a$, notre hypothèse donne
$x =f(f^{-1}(x)) = g(f^{-1}(x))$. Donc $x \lessdot g(a)$. De même, si
$x \lessdot g(a)$, alors $x \lessdot f(a)$.

En général, $(\forall x \in B)(x \lessdot f(a) \liff x \lessdot
g(a))$. Il s'ensuit que $f(a) = g(a)$ par la
\olref[sfr][rel][ord]{prop:extensionality-strictlinearorders}. Ainsi,
$(\forall a \in A)f(a) = g(a)$ par la \olref[wo]{propwoinduction}.
\end{proof}\noindent
Voilà qui donne une idée de la robustesse des bons ordres. Pour aller
plus loin, introduisons quelques notations supplémentaires.

\begin{defn}
Si $\tuple{A, <}$ est un bon ordre et si $a \in A$, posons
$A_a = \Setabs{x \in A}{x < a}$. Nous disons que $A_a$ est un
\emph{segment initial} propre de $A$ (et admettons que $A$ lui-même
soit un segment initial impropre de $A$). Soit $<_a$ la restriction de
$<$ à ce segment initial, c'est-à-dire
$\funrestrictionto{\mathord{<}}{A_a^2}$.
\end{defn}
\noindent
Avec cette notation, nous pouvons énoncer et démontrer qu'aucun bon
ordre n'est isomorphe à l'un de ses segments initiaux propres.

\begin{lem}\ollabel{wellordnotinitial}
Si $\tuple{A, <}$ est un bon ordre et si $a \in A$, alors
$\ordneq{\tuple{A, <}}{\tuple{A_a, <_a}}$.
\end{lem}

\begin{proof}
Par l'absurde, supposons que $f \colon A \to A_a$ soit un
isomorphisme. Puisque $f$ est une bijection et que $A_a \subsetneq A$,
\olref[wo]{wo:strictorder} nous permet de choisir le plus petit
!!{element} $b \in A$ pour $<$ tel que $b \neq f(b)$. Nous allons
montrer que $(\forall x \in A)(x<b \liff x < f(b))$ ; il s'ensuivra,
par la \olref[sfr][rel][ord]{prop:extensionality-strictlinearorders}, que
$b = f(b)$, ce qui donnera la contradiction cherchée.

Supposons $x < b$. Alors $x = f(x)$ par le choix de $b$. Comme $f$ est
un isomorphisme, $f(x) < f(b)$. Donc $x < f(b)$.

Supposons $x < f(b)$. Alors $f^{-1}(x) < b$, puisque $f$ est un
isomorphisme ; et donc $f^{-1}(x) = x$ par le choix de $b$. Ainsi,
$x < b$.
\end{proof}

Le résultat suivant montre, en gros, que la restriction d'un
isomorphisme à un «~segment initial~» est encore un isomorphisme :

\begin{lem}\ollabel{wellordinitialsegment}
Soient $\tuple{A, <}$ et $\tuple{B, \lessdot}$ deux bons ordres. Si
$f \colon A \to B$ est un isomorphisme et si $a \in A$, alors
$\funrestrictionto{f}{A_{a}} : A_a \to B_{f(a)}$ est un isomorphisme.
\end{lem}

\begin{proof}
Puisque $f$ est un isomorphisme :
	%Puisque f est un isomorphisme, b < a ssi f(b) \lessdot f(a), donc
\begin{align*}
	\funimage{f}{A_a} &= \funimage{f}{\Setabs{x \in A}{x < a}}\\
	&= \funimage{f}{\Setabs{f^{-1}(y) \in A}{f^{-1}(y) < a}} \\
	&= \Setabs{y \in B}{y \lessdot f(a)} \\
	&=B_{f(a)}
\end{align*}
Et $\funrestrictionto{f}{A_a}$ préserve l'ordre, puisque $f$ le
préserve.
\end{proof}

Les deux résultats qui suivent établissent que deux bons ordres sont
toujours comparables :

\begin{lem}\ollabel{lemordsegments}
Soient $\tuple{A, <}$ et $\tuple{B, \lessdot}$ deux bons ordres. Si
$\ordeq{\tuple{A_{a_1}, <_{a_1}}}{\tuple{B_{b_1}, \lessdot_{b_1}}}$
et $\ordeq{\tuple{A_{{a_2}}, <_{a_2}}}{\tuple{B_{{b_2}},
\lessdot_{b_2}}}$, alors ${a_1} < {a_2} \text{ ssi }{b_1} \lessdot
{b_2}$.
\end{lem}

\begin{proof}
Nous démontrons l'implication \emph{de gauche à droite} ; l'autre sens
est analogue. Supposons que
$\ordeq{\tuple{A_{a_1}, <_{a_1}}}{\tuple{B_{b_1},
\lessdot_{b_1}}}$ et
$\ordeq{\tuple{A_{{a_2}}, <_{a_2}}}{\tuple{B_{{b_2}},
\lessdot_{b_2}}}$, et soit $f \colon A_{{a_2}} \to B_{{b_2}}$
l'isomorphisme considéré. Supposons ${a_1} < {a_2}$ ; alors
$\ordeq{\tuple{A_{a_1}, <_{a_1}}}{\tuple{B_{f({a_1})},
\lessdot_{f({a_1})}}}$ par le \olref{wellordinitialsegment}. Ainsi,
$\ordeq{\tuple{B_{b_1}, \lessdot_{b_1}}}{\tuple{B_{f({a_1})},
\lessdot_{f({a_1})}}}$, d'où ${b_1} = f({a_1})$ par le
\olref{wellordnotinitial}. Enfin, ${b_1} \lessdot {b_2}$ puisque
l'image de $f$ est $B_{{b_2}}$.\footnote{Note de l'édition
française : l'original indique ici le domaine de $f$, alors que
$f \colon A_{a_2} \to B_{b_2}$ a pour image $B_{b_2}$.}
\end{proof}

\begin{thm}\ollabel{thm:woalwayscomparable}
Quels que soient deux bons ordres, l'un est isomorphe à un segment
initial, éventuellement impropre, de l'autre.
\end{thm}

\begin{proof}
Soient $\tuple{A, <}$ et $\tuple{B, \lessdot}$ deux bons ordres.
Par séparation, posons
\[
	f = \Setabs{\tuple{a, b} \in A \times B}{
		\ordeq{\tuple{A_a, <_a}}{\tuple{B_b, \lessdot_b}}}.
\]
Par le \olref{lemordsegments}, $a_1 < a_2$ si et seulement si
$b_1 \lessdot b_2$ pour tous
$\tuple{a_1, b_1}, \tuple{a_2, b_2} \in f$. Ainsi,
$f \colon \dom{f} \to \ran{f}$ est un isomorphisme.

Si $a_2 \in \dom{f}$ et $a_1 < a_2$, alors $a_1 \in \dom{f}$ par le
\olref{wellordinitialsegment} ; donc $\dom{f}$ est un segment initial
de $A$. De même, $\ran{f}$ est un segment initial de $B$. Par
l'absurde, supposons que les deux soient \emph{propres}. Soit alors
$a$ le plus petit !!{element} de $A \setminus \dom{f}$ pour $<$, de
sorte que $\dom{f} = A_a$, et soit $b$ le plus petit !!{element} de
$B \setminus \ran{f}$ pour $\lessdot$, de sorte que $\ran{f} = B_b$.
Ainsi, $f \colon A_a \to B_b$ est un isomorphisme, donc
$\tuple{a, b} \in f$, contradiction.
\end{proof}

\end{document}
